\begin{figure}[ht]
\centering
\begin{tikzpicture}[x=1cm, y=1cm, font=\small]

% Sum of two close sinusoids: f1 = 1.0, f2 = 1.15 (arbitrary units).
% The sum shows a fast carrier at the mean frequency modulated by a
% slow envelope at half the difference frequency.
\draw[->, thick] (-0.1, 0) -- (10.4, 0) node[right, font=\footnotesize] {$t$};
\draw[->, thick] (0, -1.4) -- (0, 1.5);

% Envelope: 2*cos(2*pi*(delta/2)*t), delta = 0.15, half-difference 0.075
\draw[domain=0:10, samples=300, smooth, engred, thick, dashed]
  plot (\x, {2*cos(deg(2*pi*0.075*\x))*0.5});
\draw[domain=0:10, samples=300, smooth, engred, thick, dashed]
  plot (\x, {-2*cos(deg(2*pi*0.075*\x))*0.5});

% The summed waveform: cos(2*pi*1.0*t) + cos(2*pi*1.15*t), scaled by 0.5
\draw[domain=0:10, samples=600, smooth, engblue, thick]
  plot (\x, {0.5*(cos(deg(2*pi*1.0*\x)) + cos(deg(2*pi*1.15*\x)))});

\node[engblue, anchor=west, font=\footnotesize] at (10.4, 0.9)
  {sum of two tones};
\node[engred, anchor=west, font=\footnotesize] at (10.4, 0.4)
  {beat envelope};

% Mark one beat period (envelope period = 1/delta ~ 6.67 in these units)
\draw[<->, enggray, thin] (0, -1.25) -- (6.67, -1.25)
  node[midway, below, font=\footnotesize] {one beat period $1/\Delta f$};

\end{tikzpicture}
\caption[Beat note from two close frequencies]{Two sinusoids of
nearly equal frequency, $f_{1}$ and $f_{2}$, sum to a fast carrier
at the mean frequency $(f_{1}+f_{2})/2$ modulated by a slow
envelope (red dashed) whose amplitude peaks recur at the beat
frequency $\Delta f = |f_{1}-f_{2}|$. Counting envelope peaks over
a known interval measures $\Delta f$ directly, the working
principle behind heterodyne frequency comparison and the audible
beats a musician uses to tune two strings to a common pitch.}
\label{fig:vol01:ch06:beat-note}
\end{figure}
